% !TeX root = ../../Book3.tex
% 纲要标题：### 6.6 导子理想与曲线正规化的比较
% 纲要位置：计划纲要.md，第 1679 行。

\section{导子理想与曲线正规化的比较}
\label{b3:sec:0606}

% TODO: 按以下纲要要点撰写本节。

% 纲要内容（仅作写作计划，尚非教材正文）：
%
% * 对有限双有理扩张 \(A\subseteq B\) 定义导子理想 \(\mathfrak c=\{a\in A:aB\subseteq A\}\)，说明其同时为 \(A\) 和 \(B\) 的理想
% * 在 \(A=k[t^2,t^3]\subseteq B=k[t]\) 中计算 \(\mathfrak c=t^2k[t]\)、\(B/A\) 与奇点处的差异
% * 以特征不为2的曲线 \(y^2=x^2(x+1)\) 为结点模型，使用 \(x=t^2-1\)、\(y=t(t^2-1)\) 核验正规化及奇点上的两个原像
% * 比较有限双有理映射、点集上的行为和局部环的改变；平坦性与非分歧性的独立检验留到第23.4节
%
% ---
%

\subsection{有限双有理扩张中的导子理想}
\label{b3:subsec:0606:01}

待撰写。

\subsection{尖点的导子理想与正规化商模}
\label{b3:subsec:0606:02}

待撰写。

\subsection{结点正规化与奇点上的两个原像}
\label{b3:subsec:0606:03}

待撰写。

\subsection{有限双有理映射的点集与局部环}
\label{b3:subsec:0606:04}

待撰写。
